% Copyright 2026 Alan Szmyt. All rights reserved pending publication license.
\section{Results}
\label{sec:results}

This design/protocol manuscript contains no formal human results. No formal
study has been run, no qualifying real-model technical evaluation has been
executed, and the originating catharsis experience remains a hypothesis-
generating observation rather than a study result. The absence of a qualifying
record is reported as \emph{not run}, \emph{not collected}, \emph{blocked}, or
\emph{not applicable}; it is never entered as a zero.

Results reporting is governed by \texttt{ANT-REPORT-RESULTS-001}, version~1.0.0.
The machine-readable contract binds each future result slot to a named source
record, the byte-locked \texttt{ANT-PROT-FEAS-001} protocol version, a frozen
analysis pointer, an evidence class, and a claim-promotion rule. It currently
registers no T0, T1, or H1 result package. This structure adopts the complete-
accounting orientation of N-of-1 and pilot-feasibility reporting guidance---
including participant flow, missingness, harms, analysis, and progression---
without describing Antidote as a completed trial or effectiveness study
\cite{vohra2015cent,eldridge2016feasibility}.

\subsection{Technical verification reporting}
\label{sec:results-technical-verification}

The D0 design record is available for inspection: the protocol identity,
content lock, reporting slots, schema mappings, equation references, and
activation gates can be deterministically checked. This is evidence that the
design artifacts exist and reconcile; it is not an execution result and does
not establish implementation correctness, feasibility, privacy, safety, or
benefit.

The repository also contains an executable synthetic session path and its
component tests, as bounded by ANT-OBS-002. Those artifacts remain development
and conformance evidence. The ANT-Q03 exit gate in issue~\#18 has not produced
the reserved \texttt{ANT-REC-T0-001} result package, so no T0 rate, count,
latency, or success claim is promoted here. Closing an issue or observing a
green workflow would not be sufficient by itself. Promotion requires the
reviewed source revision, execution environment, exact commands, planned
denominators, every terminal state, injected failure and recovery, artifact
hashes, exclusions, and an explicit claim-ledger decision. Any accepted T0
statement must retain the labels \emph{synthetic} and \emph{technical} and
cannot imply a human exposure or response.

The interface evidence figure therefore remains a visible layout reservation.
Replacing it requires a reproducible reviewed build, capture metadata, an
explicit synthetic-or-approved-data label, and redaction review. Interface
presence could demonstrate that a control is displayed; it could not
demonstrate that the underlying authority is complete or that an experience is
helpful.

\AntidoteFigure{mvp-session-evidence-interface}

\subsection{Control adherence and continuity}
\label{sec:results-control-adherence}

T1 remains \texttt{blocked-no-real-model}. No model, adapter, analyzer,
hardware profile, technical threshold set, or qualifying run manifest has been
frozen, and \texttt{ANT-REC-T1-001} does not exist. Consequently, the governed
table below contains source and analysis identities plus current evidence
states, but no requested controls, measured features, adherence values,
latencies, buffer margins, boundary distances, failure rates, or uncertainty
estimates.

\AntidoteTable{control-adherence-results}

If T1 is later activated, reporting will show all 39 planned attempts or name
every unsupported or missing attempt. Requested-versus-measured differences
will be paired by factor and seed; generation latency and verified buffer margin
will name the hardware context; semantic and waveform boundary components will
remain separate; and partial output, timeout, cancellation, failure, and
recovery will stay in their denominators. Unsupported semantic controls will be
reported as unsupported rather than assigned an invented acoustic measure. A
technically accurate tempo, sufficient buffer, or small waveform boundary does
not establish perceived continuity, affective fit, safety, or therapeutic
value.

The benchmark figure remains data-free until the T1 source package and analysis
outputs pass promotion review. Its future panels must show denominators,
uncertainty, missing runs, failures, model and hardware revisions, and links to
structured records rather than only favorable traces.

\AntidoteFigure{technical-benchmark-panels}

\subsection{Exposure and response completeness}
\label{sec:results-exposure-completeness}

H1 remains \texttt{blocked-no-collection-authority}; none of its activation
gates is satisfied and no assigned attempt, generated study artifact, exposure,
or response record is registered. The table shell names the future flow,
immediate-response, and later-response packages and analyses. Its blocked cells
are not denominators and must not be read as zero events.

\AntidoteTable{exposure-response-completeness}

After separately authorized collection, reporting would begin with every
assigned attempt and retain the transitions to attempted generation, verified
artifact, started playback, completed exposure, immediate window, and later
window. Conditions G, S, and P would be shown separately and together. Generated
but unplayed audio would remain outside the exposure denominator; an interrupted
session would remain in the attempt ledger; a replacement would not erase the
original assignment; and withdrawal would not be converted into technical
failure. Each response field would be classified as present, declined, missed-
window, interrupted, technically unavailable, or not applicable.

\subsection{Subjective response and aftereffect}
\label{sec:results-subjective-response}

No perceived-expression, felt-state, helpfulness, resonance, mismatch, harm,
surprise, burden, or aftereffect observation is reported. The reporting contract
reserves three distinct H1 views: perceived musical expression versus felt
state; immediate appraisal by condition; and immediate versus later aftereffect.
The first preserves both reports without declaring either one correct. The
second shows every session point and the prespecified within-block P-minus-S,
P-minus-G, and S-minus-G contrasts. The third retains the immediate and later
records as separate timestamped accounts rather than using later meaning to
rewrite the immediate response.

Future summaries must display exact denominators, medians, ranges, blockwise
contrasts, and all session points. The within-block permutation distribution is
eligible only if all six blocks satisfy its frozen analysis conditions and
would remain exploratory within-person uncertainty, not a population estimate
or confirmatory test. No clinical outcome is collected, and movement toward a
desired state cannot substitute for the person's appraisal of whether an
experience was wanted, helpful, mismatched, burdensome, or harmful.

Any populated report must also retain the prespecified sensitivity views:
started exposures versus those with at least 80 percent of audio heard, with
versus without carryover flags, early versus later blocks, leave-one-block-out
summaries, fixed-seed versus unsupported-seed runs, administration status, and
the declared best- and worst-case missing-value bounds. These are reporting
views rather than opportunities to select a preferred estimate after observing
the data. None has been computed.

\subsection{Null, negative, missing, and adverse observations}
\label{sec:results-nonpositive-observations}

The empty-state table gives null or opposite-direction response, mismatch or
unwanted intensity, burden, harm or adverse response, missingness, and
interruption the same structural visibility as a future favorable observation.
It contains no event counts because H1 has not begun.

\AntidoteTable{outcomes-accounting}

In a qualifying report, zero movement would remain distinct from a missing
measurement; opposite-direction movement would remain distinct from harm;
mismatch would remain distinct from unwanted intensity; and a declined answer
would remain distinct from a missed window or technical failure. Stops,
interruptions, hearing discomfort, distress, consent or privacy failures,
technical safety failures, and participant-described adverse responses would
be listed with their denominators and review dispositions. None could be
silently removed because it weakens a narrative. Conversely, the absence of an
observed adverse event in a bounded series could not establish clinical safety.

\subsection{Optional physiological observations}
\label{sec:results-optional-physiology}

Physiology is disabled under protocol version~1.0.0. No physiological source
record, analysis, panel, or result row is authorized, and no vital sign or
sensor estimate may steer generation within this protocol. A future module
would require a protocol amendment, separate consent, device and clock records,
signal-quality and artifact handling, a frozen analysis, and independent
privacy and safety review. Even then, physiology would remain supplemental: it
could not diagnose emotion or need, overwrite self-report, demonstrate
entrainment, or establish benefit.

\subsection{Promotion and replacement rule}
\label{sec:results-promotion-rule}

The current Results state is itself reproducible: its three table shells are
generated from the reporting contract, while the two contingent figures remain
visibly provisional. To replace an empty state, a future change must add the
reserved source package, validate its protocol and analysis identity, preserve
complete denominators and non-positive records, register any deviations, and
record a claim-ledger promotion decision. The reporting contract and generated
projection must then be versioned together. A source package may remain
available in the repository while its manuscript claim is rejected or
downgraded.

This gate is deliberately asymmetric: qualifying evidence may populate only
the slot it supports, while failure at any provenance, authority, analysis, or
claim boundary leaves the slot unavailable. Technical verification cannot
unlock human-response claims, H1 feasibility cannot establish efficacy, and no
result at these stages can establish a neurological mechanism.
